\documentclass[10pt,a4paper]{article}
\usepackage[utf8]{inputenc}
\usepackage[T1]{fontenc}
\usepackage[spanish]{babel}
\usepackage{amsmath}
\usepackage{amsfonts}
\usepackage{amssymb}
\usepackage{makeidx}
\usepackage{graphicx}
\usepackage{parskip}
\author{Jonathan G Araya}
\begin{document}

	\begin{center}
		\LARGE Instituto Tecnológico de Costa Rica \\
		\LARGE Área Académica Ingeniería en Computadores \\
		\LARGE CE-3201: Análisis Numérico para Ingeniería \\
	\end{center}

	\begin{figure}[h]
		\centering
		\includegraphics[scale = 0.5]{img/TEC.png}
	\end{figure}

	\begin{center}
		\huge \textbf{Manual Funtras} \\
		\vspace{20mm}
		\LARGE Integrantes \\
		\vspace{5mm}
		\large Josué Araya García \\
		\large Jonathan Guzmán Araya \\	
		\large Mariano Muñoz Masís \\
		\large Daniel Prieto Sibaja \\
		\vspace{5cm}
		Cartago, Costa Rica \\
		27 Marzo, 2021
	\end{center}

	\tableofcontents
	
	\section{Introducción}
	
	En las matemáticas existen diversos tipos de funciones como lo pueden ser:
	\begin{itemize}
		\item Algebraicas
		\item Trascendentes
	\end{itemize}

	Para este desarrollo nos enfocaremos en las funciones trascendentes, estas son  las funciones que no satisfacen una ecuación polinomial cuyos coeficientes sean a su vez polinomios; esto contrasta con las funciones algebraicas, las cuales satisfacen dicha ecuación.
	
	\subsection{¿Qué es Funtras?}
	
	Funtras es una biblioteca de funciones trascendentes desarrolladas en el lenguaje $C^{++}$ con el
	objetivo de aproximar dichas funciones mediante el uso de métodos iterativos utilizando únicamente operaciones de suma, resta, multiplicación y potencia con una cantidad de iteraciones máximas de 2500 y una tolerancia de $10^{-8}$.
		
	\section{Requisitos e Instalación}
	
	En esta sección se abarcarán los requisitos mínimos para su ejecución así como una breve guía de instalación de la misma.
	
	\subsection{Requisitos}
	El desarrollo y las pruebas de esta biblioteca se realizaron en el SO Windows 10, por lo tanto como requisitos se tiene:
	
	\begin{itemize}
		\item Sistema operativo Windows 10 o Linux
		\item MinGW
		\item Editor de código
	\end{itemize}
	
	\subsection{Instalación}
	
	\begin{itemize}
		\item Crear una carpeta donde se realizará un proyecto nuevo.
		\item Agregar a esta carpeta los archivos funtras.a y funtras.h.
		\item Crear un archivo .cpp donde se utilizará la biblioteca.
		\item En el archivo .cpp anterior aregar el header de la librería a utilizar, de la siguiente manera: include "funtras.h".
		\item Una vez realizado esto se puede utilizar la biblioteca para prograr alguna función deseada.
		\item Una vez esto se encuentre listo se deben ejecutar los siguientes comandos:
		\begin{itemize}
			\item g++ -c archivo.cpp
			\item g++ -o archivo archivo.o funtras.a
			\item archivo
		\end{itemize}
		\item Finalmente se mostrara en la terminal el resultado de la operacion utilizada en el archivo .cpp.
	\end{itemize}
		
		
	\section{Funciones implementadas en Funtras}
	A continuación se detallan las funciones implementadas en la biblioteca funtras.
	
	\subsection{Inverso multiplicativo $a^{-1}$}
	
	Esta función calcula el inverso multiplicativo, recíproco o inverso de un número $x$ real positivo, el uso de la misma se realiza de la siguiente manera:
	
	\begin{center}
			\framebox[4cm][c]{varM1(a)}
	\end{center}
	
	\subsubsection{Formulación matemática}
	
	El cálculo se realiza mediante el método iterativo que se describe a continuación.
	
	\begin{equation}\label{key1}
		x_{k+1} = x_{k}(2 - a\cdot x_{k})
	\end{equation}
		
	\subsubsection{Valores iniciales}
	
	El valor de $x_{0}$ esta dado por:
	
	\begin{align*}
		x_{0} = 
		\begin{cases}
			eps^{15}~ si~  80! < a \leq 100! \\
			eps^{11}~ si~  60! < a \leq 80! \\
			eps^{8}~~ si~  40! < a \leq 60! \\
			eps^{4}~~ si~  20! < a \leq 40! \\
			eps^{2}~~ si~~  0! < a \leq 20!
		\end{cases}
	\end{align*}
	
	donde eps es una constante ya definida con valor de:
	
	\begin{equation}\label{key2}
		eps = 2.2204x10^{-16}
	\end{equation}
	
	\subsubsection{Condición de parada}
	
	La condición de parada de la iteración está dada por: 
	
	\begin{equation}\label{key3}
		\left\lvert{\frac{x_{k+1} - x_{k}}{x_{k+1}}}\right\lvert
	\end{equation}

	Cuando la tolerancia dada sea mayor que esta, entonces devuelve el resultado obtenido.
	
	\subsubsection{Ejemplos numérico}

	double expe = exp\char`_t(11); \newline
	cout $<<$ $"$ El resultado de $e^{11}$ es: " $<<$ expe; \newline
	El resultado de $e^{11}$ es: 59874.1 \newline
	
	\subsection{Exponencial de Euler $e^{x}$}
	
	Esta función calcula el exponencial de $e$ elevado a un número natural $x$, el uso de la misma se realiza de la siguiente manera:
	
	\begin{center}
		\framebox[4cm][c]{exp{\_}t(x)}
	\end{center}
	
	\subsubsection{Formulación matemática}
	
	El cálculo se realiza mediante la sumatoria que se describe a continuación.
	
	\begin{equation}\label{key4}
		S_{k}(a) = \sum_{n=0}^{k}\frac{a^{n}}{n!}
	\end{equation}
	
	\subsubsection{Condición de parada}
	
	La condición de parada de la iteración está dada por: 
	
	\begin{equation}\label{key5}
		\left\lvert S_{k+1}(a) - S_{k}(a) \right\lvert < tol
	\end{equation}
	
	\subsubsection{Ejemplo numérico}

	double expoNUNO = varM1(144); \newline
	cout $<<$ $"$ El resultado de $144^{-1}$ es: " $<<$ expoNUNO; \newline
	El resultado de $144^{-1}$ es: 0.00694444\newline

	\subsection{Seno $\sin(x)$}
	
	Esta función calcula el $seno$ de un número $x$, el uso de la misma se realiza de la siguiente manera:
	
	\begin{center}
		\framebox[4cm][c]{sin{\_}t(x)}
	\end{center}
	
	\subsubsection{Formulación matemática}
	
	El cálculo se realiza mediante la sumatoria que se describe a continuación.
	
	\begin{equation}\label{key6}
		S_{k}(a) = \sum_{n=0}^{k}(-1)^{n}\frac{a^{2n + 1}}{(2n + 1)!}
	\end{equation}
	
	\subsubsection{Condición de parada}
	
	La condición de parada de la iteración está dada por: 
	
	\begin{equation}\label{key7}
		\left\lvert S_{k+1}(a) - S_{k}(a) \right\lvert < tol
	\end{equation}
	
	
	\subsubsection{Ejemplo numérico}
	
	double seno = sin{\_}t(2.33); \newline
	cout $<<$ $"$ El resultado de $sen(2.33)$ es: " $<<$ seno; \newline
	El resultado de $sen(2.33)$ es: 0.725384\newline

	\subsection{Coseno $\cos(x)$}
	
	Esta función calcula el $coseno$ de un número $x$, el uso de la misma se realiza de la siguiente manera:
	
		\begin{center}
		\framebox[4cm][c]{cos{\_}t(x)}
	\end{center}
	
	\subsubsection{Formulación matemática}
	
	El cálculo se realiza mediante la sumatoria que se describe a continuación.
	
	\begin{equation}\label{key8}
		S_{k}(a) = \sum_{n=0}^{k}(-1)^{n}\frac{a^{2n}}{(2n)!}
	\end{equation}
	
	\subsubsection{Condición de parada}
	
	La condición de parada de la iteración está dada por: 
	
	\begin{equation}\label{key9}
		\left\lvert S_{k+1}(a) - S_{k}(a) \right\lvert < tol
	\end{equation}
	
	\subsubsection{Ejemplo numérico}

	double coseno = cos{\_}t(3.78); \newline
	cout $<<$ $"$ El resultado de $cos(3.78)$ es: " $<<$ coseno; \newline
	El resultado de $cos(3.78)$ es: -0.803046\newline

	\subsection{Tangente $\tan(x)$}
	
	Esta función calcula la $tangente$ de un número $x$, el uso de la misma se realiza de la siguiente manera:
	
	\begin{center}
		\framebox[4cm][c]{tan{\_}t(x)}
	\end{center}
	
	\subsubsection{Formulación matemática}
	
	La función tangente se puede componer a partir de otras como lo son $seno$ y $coseno$, es por ello que el calculo de la misma se realiza mediante la siguiente ecuación:
	
	\begin{equation}\label{key10}
		tan(x) = sen(x)\cdot cos(x)^{-1}
	\end{equation}
	
	\subsubsection{Ejemplo numérico}

	double tangente = tan{\_}t(7); \newline
	cout $<<$ $"$ El resultado de $tan(7)$ es: " $<<$ tangente; \newline
	El resultado de $tan(7)$ es: 0.871449\newline
	
	\subsection{Logaritmo natural $\ln(x)$}
	
	Esta función calcula el $logaritmo natural$ de un número $x$, el uso de la misma se realiza de la siguiente manera:
	
	\begin{center}
		\framebox[4cm][c]{ln{\_}t(x)}
	\end{center}
	
	\subsubsection{Formulación matemática}
	
	El cálculo se realiza mediante la sumatoria que se describe a continuación.
	
	\begin{equation}\label{key11}
		S_{k}(a) = \frac{2(a-1)}{a + 1}\sum_{n=0}^{k}\frac{1}{2n + 1}\left(\frac{a - 1}{a + 1}\right)^{2n}
	\end{equation}
	
	\subsubsection{Condición de parada}
	
	La condición de parada de la iteración está dada por: 
	
	\begin{equation}\label{key12}
		\left\lvert S_{k+1}(a) - S_{k}(a) \right\lvert < tol
	\end{equation}
	
	\subsubsection{Ejemplo numérico}

	double logan = ln{\_}t(45); \newline
	cout $<<$ $"$ El resultado de $ln(45)$ es: " $<<$ logan; \newline
	El resultado de $ln(45)$ es:3.80666\newline
	
	\subsection{Logaritmo $\log_{a}(x)$}
	
	Esta función calcula el $logaritmo$ de base $a$ a un número $x$, el uso de la misma se realiza de la siguiente manera:
	
	\begin{center}
		\framebox[4cm][c]{log{\_}t(x)}
	\end{center}
	
	\subsubsection{Formulación matemática}
	
	La función logaritmo se puede componer a partir de otra como $logaritmo natural$, es por ello que el calculo de la misma se realiza mediante la siguiente ecuación:
	
	\begin{equation}\label{key13}
		\log_{a}(x) = \ln(x)\cdot (\ln(a))^{-1}
	\end{equation}
	
	\subsubsection{Ejemplo numérico}

	double logb = log{\_}t(33, 43); \newline
	cout $<<$ $"$ El resultado de logaritmo en base 33 de 43 es: es: " $<<$ logb;
	El resultado de  logaritmo en base 33 de 43 es: 1.0757\newline
	
	\subsection{Exponencial $a^{x}$}
	
	Esta función calcula el exponencial de un número $a$ elevado a un número $x$, el uso de la misma se realiza de la siguiente manera:
	
	\begin{center}
		\framebox[4cm][c]{power{\_}t(x)}
	\end{center}
	
	\subsubsection{Formulación matemática}
	Esta función se compone de funciones ya programadas anteriormente, de esta forma podemos ver la siguiente relación:
	\begin{equation}\label{key13}
		a^{x} = e^{a \cdot \ln(x)}
	\end{equation}
	Por lo tanto para esta función se aplican la función de logaritmo natural y exponencial de euler anteriormente explicadas.
	
	\subsubsection{Valores iniciales}
	
	Está función no tiene valores iniciales, solo los dependientes a las funciones que la componen.		

	\subsubsection{Condición de parada}

	Esta función no tiene condición de parada como tal, solo las dependientes a las funciones que la componen.
	
	\subsubsection{Ejemplo numérico}

	double powert = power{\_}t(2, 11); \newline
	cout $<<$ $"$ El resultado de 2 elevado a la 11 es: " $<<$ powert; \newline
	El resultado de 2 elevado a la 11 es: 2048\newline
	
	\subsection{Seno hiperbólico $\sinh(x)$}
	
	Esta función calcula el $seno hiperbólico$ de un número $x$, el uso de la misma se realiza de la siguiente manera:
	
	\begin{center}
		\framebox[4cm][c]{sinh{\_}t(x)}
	\end{center}
	
	\subsubsection{Formulación matemática}
	
	El cálculo se realiza mediante la sumatoria que se describe a continuación.
	
	\begin{equation}\label{key14}
		S_{k}(a) = \sum_{n=0}^{k}\frac{a^{2n + 1}}{(2n + 1)!}
	\end{equation}
	
	\subsubsection{Condición de parada}
	
	La condición de parada de la iteración está dada por: 
	
	\begin{equation}\label{key15}
		\left\lvert S_{k+1}(a) - S_{k}(a) \right\lvert < tol
	\end{equation}
		
	\subsubsection{Ejemplo numérico}

	double senoh = sinh{\_}t(2.1); \newline
	cout $<<$ $"$ El resultado de $sinh(2.1)$ es: " $<<$ senoh; \newline
	El resultado de $sinh(2.1)$ es: 4.02186\newline
	
	\subsection{Coseno hiperbólico $\cosh(x)$}
	
	Esta función calcula el $coseno hiperbólico$ de un número $x$, el uso de la misma se realiza de la siguiente manera:
	
	\begin{center}
		\framebox[4cm][c]{cosh{\_}t(x)}
	\end{center}
	
	\subsubsection{Formulación matemática}

	El cálculo se realiza mediante la sumatoria que se describe a continuación.

	\begin{equation}\label{key16}
		S_{k}(a) = \sum_{n=0}^{k}\frac{a^{2n}}{(2n)!}
	\end{equation}

	\subsubsection{Condición de parada}

	La condición de parada de la iteración está dada por: 

	\begin{equation}\label{key17}
		\left\lvert S_{k+1}(a) - S_{k}(a) \right\lvert < tol
	\end{equation}
	
	\subsubsection{Ejemplo numérico}
	
	double cosenoh = cosh{\_}t(0.89); \newline
	cout $<<$ $"$ El resultado de $cosh(0.89)$ es: " $<<$ cosenoh; \newline
	El resultado de $cosh(0.89)$ es: 1.42289\newline

	\subsection{Tangente hiperbólico $\tanh(x)$}
	
	Esta función calcula la $tangente hiperbólica$ de un número $x$, el uso de la misma se realiza de la siguiente manera:
	
	\begin{center}
		\framebox[4cm][c]{tanh{\_}t(x)}
	\end{center}
	
	\subsubsection{Formulación matemática}
	
	La función tangente hiperbólico se puede componer a partir de otras como lo son $seno hiperbólico$ y $coseno hiperbólico$, es por ello que el calculo de la misma se realiza mediante la siguiente ecuación:
	
	\begin{equation}\label{key18}
		tanh(x) = senh(x)\cdot cosh(x)^{-1}
	\end{equation}
	
	\subsubsection{Ejemplo numérico}

	double tanoh = tanh{\_}t(0.22); \newline
	cout $<<$ $"$ El resultado de $tanh(0.22)$ es: " $<<$ tanoh; \newline
	El resultado de $tanh(0.22)$ es: 0.216518\newline
	
	\subsection{Raíz cuadrada $\sqrt{x}$}
	
	Esta función calcula la $raíz cuadrada$ de un número $x$ mediante el método de Newton-Raphson, el uso de la misma se realiza de la siguiente manera:
	
	\begin{center}
		\framebox[4cm][c]{sqrt{\_}t(x)}
	\end{center}
	
	\subsubsection{Formulación matemática}
	
	El método de Newton-Raphson se define como:
	
	\begin{equation}\label{key19}
		x_{k+1} = x_{k} - \frac{f(x_{k})}{f'(x_{k})}, ~x_{0} = \alpha
	\end{equation}
	
	Donde para encontrar la p-ésima raíz de a:
	
	\begin{equation}\label{key20}
		\sqrt[p]{a}	
	\end{equation}

	El cero de la función está dado por:
	
	\begin{equation}\label{key21}
		g(x) = x^{p} - a, ~x_{0} = \frac{a}{2}
	\end{equation}

	Por lo que resulta en la siguiente iteración:
	
	\begin{align*}
		\centering
		\begin{cases}
			x_{0} = \frac{a}{2} \\
			x_{k+1} = x_{k} - \frac{x_{k}^{2} - a}{2\cdot x_{k}}
		\end{cases}
	\end{align*}
	
	
	\subsubsection{Valore inicial}
	El valor inicial está dado por: $x_{0 = \frac{a}{2}}$.
	\subsubsection{Condición de parada}
	
	\begin{equation}\label{key22}
		\left\lvert (S_{k+1}(a) - S_{k}(a)) / S_{k+1}(a) \right\lvert < tol
	\end{equation}
	
	\subsubsection{Ejemplo numérico}
	double raiz = sqrt{\_}t(144); \newline
	cout $<<$ $"$ El resultado de la raíz cuadrada de 144 es: " $<<$ raiz; \newline
	El resultado de la raíz cuadrada de 144 es:12\newline
	
	
	\subsection{Raíz $\sqrt[a]{x}$}
	
	Esta función calcula la $raíz a-ésima$ de un número $x$ mediante el método de Newton-Raphson, el uso de la misma se realiza de la siguiente manera:
	
	\begin{center}
		\framebox[4cm][c]{root{\_}t(x)}
	\end{center}
	
	\subsubsection{Formulación matemática}

	El método de Newton-Raphson se define como:
	
	\begin{equation}\label{key23}
		x_{k+1} = x_{k} - \frac{f(x_{k})}{f'(x_{k})}, ~x_{0} = \alpha
	\end{equation}
	
	Donde para encontrar la p-ésima raíz de a:
	
	\begin{equation}\label{key24}
		\sqrt[p]{a}	
	\end{equation}

	El cero de la función está dado por:
	
	\begin{equation}\label{key25}
		g(x) = x^{p} - a, ~x_{0} = \frac{a}{2}
	\end{equation}

	Por lo que resulta en la siguiente iteración:
	
	\begin{align*}
		\centering
		\begin{cases}
			x_{0} = \frac{a}{2} \\
			x_{k+1} = x_{k} - \frac{x_{k}^{p} - a}{2\cdot x_{k}}
		\end{cases}
	\end{align*}
	
	\subsubsection{Valores iniciales}
	El valor inicial está dado por: $x_{0 = \frac{a}{2}}$.
	\subsubsection{Condición de parada}
	\begin{equation}\label{key22}
		\left\lvert (S_{k+1}(a) - S_{k}(a)) / S_{k+1}(a) \right\lvert < tol
	\end{equation}
	\subsubsection{Ejemplo numérico}
	
	double raizx = root{\_}t(45, 7); \newline
	cout $<<$ $"$ El resultado de la raíz 7 de 45 es: " $<<$ raizx; \newline
	El resultado de la raíz 7 de 45 es:1.72256\newline
	
	\subsection{Arseno $\sin^{-1}(x)$}
	
	Esta función calcula el $arcoseno$ de un número $x$, el uso de la misma se realiza de la siguiente manera:
	
	\begin{center}
		\framebox[4cm][c]{asin{\_}t(x)}
	\end{center}
	
	\subsubsection{Formulación matemática}
	
	El cálculo se realiza mediante la sumatoria que se describe a continuación.
	
	\begin{equation}\label{key22}
		S_{k}(a) = \sum_{n=0}^{k}\frac{2n!}{4^{n}(n!)^{2}(2n + 1)}a^{2n+1}
	\end{equation}
	
	\subsubsection{Condición de parada}
	
	La condición de parada de la iteración está dada por: 
	
	\begin{equation}\label{key9}
		\left\lvert S_{k+1}(a) - S_{k}(a) \right\lvert < tol
	\end{equation}
	
	\subsubsection{Ejemplo numérico}

	double aseno = asin{\_}t(0.33); \newline
	cout $<<$ $"$ El resultado de aseno(0.33) es: " $<<$ aseno; \newline
	El resultado de aseno(0.33) es: 0.336304\newline
	
	\subsection{Arcotangente $\tan^{-1}(x)$}
	
	Esta función calcula el $arcotangente$ de un número $x$, el uso de la misma se realiza de la siguiente manera:
	
	\begin{center}
		\framebox[4cm][c]{atan{\_}t(x)}
	\end{center}
	
	\subsubsection{Formulación matemática}
	
	El cálculo se realiza mediante la sumatoria que se describe a continuación.
	
	\begin{equation}\label{key22}
		S_{k}(a) = \sum_{n=0}^{k}(-1)^{n}\frac{a^{2n+1}}{2n + 1}
	\end{equation}
	
	\subsubsection{Condición de parada}
	
	La condición de parada de la iteración está dada por: 
	
	\begin{equation}\label{key9}
		\left\lvert S_{k+1}(a) - S_{k}(a) \right\lvert < tol
	\end{equation}
	
	\subsubsection{Ejemplo numérico}

	double atan = atan{\_}t(0.67); \newline
	cout $<<$ $"$ El resultado de atan(0.67) es: " $<<$ atan; \newline
	El resultado de atan(0.67) es: 0.590307\newline
	
	
\end{document}